\section{Caching sequence-search results}
\label{sec:appendix-query-cache}

The separate \texttt{bioseq\_cache} extension can save a completed search
result and let another execution read it. It can also make an execution wait
for an existing producer of the same result. This is the result cache from
Section~\ref{sec:bioseq-result-cache}. It is separate from caches of sequence
fragments, decompressed columns, and index data.

\begin{quote}
\textbf{Experimental feature.} The query cache has not been tested enough
to establish reliable result correctness and may return incorrect results.
Keep an uncached reference when evaluating it. The following examples explain
the implemented controls and intended checks, not a production guarantee.
\end{quote}

\subsection{Enable the extension and check startup}

Shared cache state is allocated when PostgreSQL starts. Installing the SQL
extension alone does not allocate this state.
\begin{enumerate}
  \item Add both libraries to \texttt{shared\_preload\_libraries}, with
    \texttt{bioseq\_cache} before \texttt{bioseq}. Retain any other libraries
    already required by the server.
  \item Restart PostgreSQL. Reloading the configuration is insufficient.
  \item Install both SQL extensions in the database, then run the startup
    checks as an administrator.
\end{enumerate}

\par\noindent\begin{minipage}{\linewidth}
\begin{lstlisting}[style=bioseq-appendix-text,
  caption={Server configuration: restart after changing this setting.}]
shared_preload_libraries = 'bioseq_cache,bioseq'
\end{lstlisting}
\end{minipage}

\par\noindent\begin{minipage}{\linewidth}
\begin{lstlisting}[style=bioseq-appendix-sql,
  caption={Install and check the extensions after restart.}]
-- Omit extensions already installed.
CREATE EXTENSION bioseq_cache;
CREATE EXTENSION bioseq;

SHOW shared_preload_libraries;
SELECT bioseq_cache.version();
SELECT count(*) FROM bioseq_cache.entries();
\end{lstlisting}
\end{minipage}

The preload list must have the required order and the entry query must
complete. A missing shared-state error, such as \texttt{requested tranche
is not registered}, calls for a corrected startup configuration and restart.
A successful \texttt{SHOW} of the mode alone is insufficient. Before loading
the library, PostgreSQL can accept its name as an unchecked custom setting.
After registration, the actual mode check rejects an enabled mode when the
required shared state was not requested at startup.

\subsection{Choose a mode and retention time}

The two cache settings are session GUCs. Their types, ranges, and defaults
also appear in Section~\ref{sec:appendix-configuration}.
\begin{itemize}
  \item \texttt{disabled}, the default, executes the retained query plan
    without query-result caching.
  \item \texttt{exact\_query} saves and replays an eligible complete
    statement result. Matching includes the analyzed request, parameters,
    settings, security context, and data dependencies. Equal SQL text alone
    does not establish compatibility.
  \item \texttt{reusable\_fragments} additionally enables the generic
    cache's coverage interface. A provider using that interface must describe
    covered input rows and identify individual result occurrences.
\end{itemize}

\textbf{The current BioSeqDB adapter supports exact results only.}
Its scan and batch-join adapters do not implement the coverage and occurrence
identity callbacks needed for partial reuse. Selecting
\texttt{reusable\_fragments} does not add these callbacks. An eligible
BioSeqDB query can still use exact-result caching in that mode. Creating an
index of type \texttt{minimizers\_per\_sequence} or
\texttt{minimizers\_per\_fragment} does not change this limitation.

\par\noindent\begin{minipage}{\linewidth}
\begin{lstlisting}[style=bioseq-appendix-sql,
  caption={Inspect the defaults and enable exact-result caching.}]
SHOW bioseq_cache.query_cache.mode;
SHOW bioseq_cache.query_cache.completed_result_ttl_ms;
SET bioseq_cache.query_cache.mode = 'exact_query';
SET bioseq_cache.query_cache.completed_result_ttl_ms = 3600000;
\end{lstlisting}
\end{minipage}

The completed-result time to live, or TTL, is measured in milliseconds.
Its default is 3\,600\,000 milliseconds, or one hour.
\begin{itemize}
  \item The retention clock starts when production completes. It is not
    a query timeout and does not stop a long-running producer.
  \item A reader uses the smaller of the producer's TTL and its own current
    TTL. Increasing the reader's setting cannot extend the producer's limit.
  \item Zero prevents later readers from reusing a completed result. A
    waiter already registered for that running producer may consume its
    newly completed result once, subject to the other validity checks.
  \item Dependency checks still apply. A long TTL does not permit reuse
    after a relevant committed change.
\end{itemize}

\subsection{Run a cold search and repeat it}

The following small table is shared by the remaining examples. Automatic
insert maintenance is enabled so the sequence writes receive the corresponding
BioSeqDB maintenance. This setting controls index maintenance, not caching.
Adjacent SQL string literals separated by a newline form one string.

\par\noindent\begin{minipage}{\linewidth}
\begin{lstlisting}[style=bioseq-appendix-sql,
  caption={Prepare a small sequence table for cache inspection.}]
SET bioseq.index.insert_maintenance.enabled = on;
CREATE TABLE cache_reads (
  pk integer PRIMARY KEY,
  seq0 bioseq.seq('dna4', 32)
) USING bioseq;
INSERT INTO cache_reads VALUES
  (1, 'TCACTGTAGGATCGTAGCTAGTCCGTACTGTAGCTGATCGTAGC'
      'CCTAGTCACTGTAGCCTGATCGTAGCTAGTCACTGTAGCCATACCTGAA'),
  (2, 'CCCCAACTTAGTCACTTCTCAAAAATTACTCTCTCATCATCATCTACT'
      'CCTAGTCACTGTAGCCTGATCGTAGCTAGTCACACACACCCCAAA'),
  (3, 'CTAGTCACTGTAGCCTGATCGTAGCTAGT');

CREATE TEMP TABLE cache_stats_before AS
SELECT * FROM bioseq_cache.stats();
\end{lstlisting}
\end{minipage}

Run the identical search twice. A \emph{cold} execution has no compatible
completed result available. A \emph{warm} execution can read one produced
earlier, if it remains valid.
The count is convenient for observing cache activity. Equal counts alone do
not establish result correctness. Compare the full result rows and alignment
fields with an uncached execution when checking correctness.

\par\noindent\begin{minipage}{\linewidth}
\begin{lstlisting}[style=bioseq-appendix-sql,
  caption={Repeat one exact request and inspect counter changes.}]
SELECT count(*) FROM cache_reads
WHERE public.mmsearch(seq0,
  'CTAGTCACTGTAGCCTGATCGTAGCTAGT'::bioseq.seq('dna4', 32));

SELECT count(*) FROM cache_reads
WHERE public.mmsearch(seq0,
  'CTAGTCACTGTAGCCTGATCGTAGCTAGT'::bioseq.seq('dna4', 32));

SELECT s.produced - b.produced AS produced_delta,
       s.completed_hits - b.completed_hits AS hit_delta,
       s.inflight_waits - b.inflight_waits AS wait_delta,
       s.reuse_hits - b.reuse_hits AS reuse_delta,
       s.fallbacks - b.fallbacks AS fallback_delta,
       s.invalidations - b.invalidations AS invalidation_delta
FROM bioseq_cache.stats() AS s
CROSS JOIN cache_stats_before AS b;
\end{lstlisting}
\end{minipage}

For an eligible cold execution followed by an exact replay, expect a completed
production followed by a completed hit. The counters cover the server, so other
sessions can contribute even to these differences. A nonzero fallback count
calls for inspecting the actual plan action and provider readiness.

\subsection{Observe one producer and its waiters}

An entry is one attempt to produce a compatible result. Its \emph{producer}
computes the result. A \emph{waiter} has found that running attempt and waits
for it. An entry's \emph{generation} distinguishes that attempt from earlier
attempts for the same request. Coordination is automatic. There is no SQL
lock or unlock call for a cache entry.
\begin{enumerate}
  \item In session A, enable exact mode and start a cold search. Use a real
    query long enough to overlap another session. The small table above will
    usually finish too quickly to demonstrate waiting.
  \item In session B, enable the same mode and submit the same statement
    while A runs. Keep parameters, options, role, and relevant settings the
    same. Avoid changing the request while testing coordination.
  \item In an administrator or monitoring session C, inspect building
    entries with the query below.
\end{enumerate}

\par\noindent\begin{minipage}{\linewidth}
\begin{lstlisting}[style=bioseq-appendix-sql,
  caption={Inspect producers and waiting executions.}]
SELECT entry_key, generation, kind, state,
       producer_pid, waiters, readers, age, ttl, sql
FROM bioseq_cache.entries()
WHERE state = 'BUILDING'
ORDER BY created_at;

SELECT produced, completed_hits, inflight_waits, fallbacks
FROM bioseq_cache.stats();
\end{lstlisting}
\end{minipage}

A short shared-state lock selects the producer. The producer then computes
without holding that lock for the whole query. Waiters sleep and check
cancellation, producer liveness, and validity. One cold producer and one
successful waiting reader should add one production, one completed hit, and
one in-flight wait to the counters.

The producer may return its own rows while recording them. Other executions
read the cached result only after completion and validation. Canceling a
waiter releases its participation and does not cancel the producer. Canceling
the producer instead follows the failure path below.

\subsection{Read the actual cache decision}

Plain \texttt{EXPLAIN} describes the planned cache wrapper.
\texttt{EXPLAIN ANALYZE} executes the statement and adds its actual action.

\par\noindent\begin{minipage}{\linewidth}
\begin{lstlisting}[style=bioseq-appendix-sql,
  caption={Inspect a planned cache route, then its execution.}]
EXPLAIN (VERBOSE, COSTS OFF)
SELECT count(*) FROM cache_reads
WHERE public.mmsearch(seq0,
  'CTAGTCACTGTAGCCTGATCGTAGCTAGT'::bioseq.seq('dna4', 32));

EXPLAIN (ANALYZE, VERBOSE, COSTS OFF)
SELECT count(*) FROM cache_reads
WHERE public.mmsearch(seq0,
  'CTAGTCACTGTAGCCTGATCGTAGCTAGT'::bioseq.seq('dna4', 32));
\end{lstlisting}
\end{minipage}

For an exact \texttt{BioseqCacheScan}, \texttt{Cache Entry Kind} is
\texttt{ExactFinalResult}. Read \texttt{Actual Cache Action} first:
\begin{itemize}
  \item \texttt{[Hit] exact final result}: read a completed exact result.
  \item \texttt{[Produce] exact final result}: produce a result for reuse.
  \item \texttt{[Bypass] retained child}: execute the underlying query plan.
\end{itemize}
Then inspect forced bypass, runtime provider validation, and dependencies.
Readiness, storage, ownership, and cleanup fields are status snapshots.
They do not establish that a particular execution returned correct rows.
Section~\ref{sec:appendix-explain} gives the complete field list and reading
order.

\subsection{Understand partial coverage and its current limit}

\emph{Coverage} describes the input rows for which a computation was carried
out. It includes rows that produced no alignment. A reusable provider must
describe coverage separately from the output rows and must distinguish
multiple alignment occurrences from the same input row.

Consider two requests with the same sequence query. The first searches input
rows with keys in $\{1,2\}$ and the second searches keys in $\{1,3\}$.
If a provider supports the generic coverage interface, the second request can:
\begin{itemize}
  \item read the completed work for input key 1;
  \item exclude input key 2 because the second request does not ask for it;
  \item compute input key 3, the \emph{residual}, or missing part.
\end{itemize}
Compatibility also depends on the provider recipe, parameters, security
context, and data versions. Overlapping SQL predicates alone are insufficient.
Section~\ref{sec:bioseq-result-cache} describes this contract formally.

The following prepared statements make the two requests explicit. They are
valid requests to inspect, but the current BioSeqDB adapter cannot combine
their overlapping coverage. It can run the retained plan or use an eligible
exact result. A repeated identical query in reusable mode is therefore not
a demonstration of partial reuse.

\par\noindent\begin{minipage}{\linewidth}
\begin{lstlisting}[style=bioseq-appendix-sql,
  caption={Two related requests; partial BioSeqDB reuse is not implemented.}]
SET bioseq_cache.query_cache.mode = 'reusable_fragments';

PREPARE cache_coverage_12(text) AS
SELECT pk FROM cache_reads
WHERE pk IN (1, 2)
  AND public.mmsearch(seq0, $1::bioseq.seq('dna4', 32));

PREPARE cache_coverage_13(text) AS
SELECT pk FROM cache_reads
WHERE pk IN (1, 3)
  AND public.mmsearch(seq0, $1::bioseq.seq('dna4', 32));

EXECUTE cache_coverage_12('CTAGTCACTGTAGCCTGATCGTAGCTAGT');
EXPLAIN (ANALYZE, VERBOSE, COSTS OFF, FORMAT JSON)
EXECUTE cache_coverage_13('CTAGTCACTGTAGCCTGATCGTAGCTAGT');

DEALLOCATE cache_coverage_12;
DEALLOCATE cache_coverage_13;
SET bioseq_cache.query_cache.mode = 'exact_query';
\end{lstlisting}
\end{minipage}

For a provider that does support coverage, \texttt{ReusableCoverage} identifies
that cache node. Its actual action distinguishes reading a complete current
generation, combining selected branches with fresh residual work, producing
a new generation, and bypassing to the retained child. To determine what
was reused, read each \texttt{Cache Branch} action and optional recheck,
then the \texttt{Residual Child} predicate. The optional \texttt{Provider}
section supplies that provider's coverage details. A
\texttt{planned/provisional} state means no current entry snapshot was
available, not that reuse occurred.

\subsection{Inspect entries and unfinished work}

\par\noindent\begin{minipage}{\linewidth}
\begin{lstlisting}[style=bioseq-appendix-sql,
  caption={Inspect individual cache generations.}]
SELECT entry_key, generation, kind, state,
       created_at, completed_at, age, ttl, rows,
       logical_bytes, stored_bytes, waiters, readers, hits,
       dependencies
FROM bioseq_cache.entries()
ORDER BY created_at DESC;

SELECT entry_key, generation, kind, state,
       producer_pid, age, waiters, readers, sql
FROM bioseq_cache.entries()
WHERE state <> 'COMPLETE'
ORDER BY created_at DESC;
\end{lstlisting}
\end{minipage}

The first query shows current entries. The second narrows this to building
or otherwise unfinished entries, including failed entries that have not yet
been removed. Interpret their values as follows:
\begin{itemize}
  \item \texttt{age} includes time spent producing the entry. The
    completed-result TTL clock starts later, at completion. Comparing
    \texttt{age} directly with \texttt{ttl} is therefore not an expiry test.
  \item \texttt{ttl} reports the producer's retention duration, not the
    remaining lifetime. The reusable EXPLAIN label \texttt{Effective TTL}
    also displays this producer value. The reader's smaller TTL can still
    restrict reuse.
  \item \texttt{logical\_bytes} is the logical result payload size and
    \texttt{stored\_bytes} is its stored payload size. Their difference
    indicates payload reduction. Neither includes all cache overhead or
    measures the PostgreSQL table size.
  \item \texttt{waiters} and \texttt{readers} describe current users.
    \texttt{hits} counts hits for exact entries. The current reusable-entry
    diagnostic leaves that field at zero, so it cannot measure reuse there.
  \item The full entry row also includes database and owner identifiers,
    producer SQL and plan text. Visibility depends on monitoring privileges.
\end{itemize}

Production records rows before publishing a completed generation. If cache
capture fails, incomplete stored rows are discarded instead of being published
as a complete result. An early stop, such as closing a cursor before consuming
the child, also cannot publish that incomplete result. Cache work that cannot
continue may fall back to the retained plan. An error in the underlying query
still raises that query error.

After a failure, inspect unfinished entries and changes in
\texttt{fallbacks}, \texttt{invalidations}, and \texttt{evictions}.
Failed, expired, or invalidated entries are removed after their users release
them. Use the administrative SQL interface for an explicit reset. Do not
delete shared cache files or alter shared state by hand.

\subsection{Check writes and administrative invalidation}

The implementation tracks relation dependencies and retires dependent cache
generations around committed writes. It is intended to prevent subsequent
requests from using a result for the old data version. Check this behavior
against an uncached query when evaluating the experimental cache.

\par\noindent\begin{minipage}{\linewidth}
\begin{lstlisting}[style=bioseq-appendix-sql,
  caption={Change an input row and inspect the next search.}]
BEGIN;
SET LOCAL bioseq.index.insert_maintenance.enabled = on;
UPDATE cache_reads
SET seq0 = 'ACGTACGTACGTACGTACGTACGTACGTACGT'
WHERE pk = 1;
COMMIT;

EXPLAIN (ANALYZE, VERBOSE, COSTS OFF)
SELECT count(*) FROM cache_reads
WHERE public.mmsearch(seq0,
  'CTAGTCACTGTAGCCTGATCGTAGCTAGT'::bioseq.seq('dna4', 32));
\end{lstlisting}
\end{minipage}

An aborted transaction does not publish its data change as a committed
invalidation. Rolling back a savepoint discards the corresponding recorded
write state. TTL is a separate retention limit and does not bypass dependency
validation.

Only a superuser can run the following administrative functions. Invalidating
entries and resetting counters are separate operations.

\par\noindent\begin{minipage}{\linewidth}
\begin{lstlisting}[style=bioseq-appendix-sql,
  caption={Administrative cache operations.}]
SELECT bioseq_cache.invalidate_relation('cache_reads'::regclass);
SELECT bioseq_cache.invalidate_all();
SELECT bioseq_cache.reset_stats();
\end{lstlisting}
\end{minipage}

An ordinary role calling \texttt{invalidate\_relation},
\texttt{invalidate\_all}, or \texttt{reset\_stats} receives
\texttt{insufficient\_privilege}, SQLSTATE \texttt{42501}.
Granting execution permission alone does not remove their superuser check.
Monitoring access does not grant administrative access.

\subsection{Read counters and grant monitoring access}

\begin{center}
\small
\begin{tabular}{@{}p{.27\textwidth}p{.68\textwidth}@{}}
\toprule
Column of \texttt{stats()} & Meaning \\
\midrule
\texttt{produced} & Completed productions. \\
\texttt{completed\_hits} & Replays of completed exact results. \\
\texttt{inflight\_waits} & Consumers that waited for an active producer. \\
\texttt{reuse\_hits} & Hits through the generic reusable-coverage route. \\
\texttt{fallbacks} & Cache routes or captures that could not be used or
  completed and retained ordinary execution. \\
\texttt{invalidations} & Invalidations recorded by the cache. \\
\texttt{evictions} & Entries removed by capacity or lifecycle policy. \\
\texttt{logical\_bytes} & Logical payload bytes produced since reset. \\
\texttt{stored\_bytes} & Stored payload bytes produced since reset. \\
\bottomrule
\end{tabular}
\end{center}

\par\noindent\begin{minipage}{\linewidth}
\begin{lstlisting}[style=bioseq-appendix-sql,
  caption={Read all global counters.}]
SELECT produced, completed_hits, inflight_waits, reuse_hits,
       fallbacks, invalidations, evictions,
       logical_bytes, stored_bytes
FROM bioseq_cache.stats();
\end{lstlisting}
\end{minipage}

The two global byte counters are cumulative production totals, not the
currently occupied cache memory or disk space. For a listed generation, use
the entry's payload sizes. The EXPLAIN \texttt{Cache Storage} diagnostic
reports reservations and ceilings for memory and file storage. The extension
does not expose a cache-capacity GUC in this interface.

Take before-and-after snapshots around a controlled workload, as in the exact
search example. Keep the test session otherwise idle and account for concurrent
traffic. Repeated fallbacks suggest checking eligibility and readiness.
Frequent invalidations suggest that writes retire results before they are
reused. For large stored results, first narrow the request and inspect its
actual cache action.

By default, \texttt{stats()} and \texttt{entries()} are accessible through
membership in PostgreSQL's \texttt{pg\_read\_all\_stats} role.
An administrator can grant that role to an existing observer:

\par\noindent\begin{minipage}{\linewidth}
\begin{lstlisting}[style=bioseq-appendix-sql,
  caption={Grant the default monitoring role.}]
GRANT pg_read_all_stats TO cache_observer;
\end{lstlisting}
\end{minipage}

Alternatively, an administrator can grant schema usage and execution of the
two monitoring functions directly. Those narrower grants allow access without
membership in \texttt{pg\_read\_all\_stats}. Producer details can remain
redacted for such an observer. A superuser, the producer's role, or a member
of \texttt{pg\_read\_all\_stats} can see the fuller diagnostic detail.
Grant monitoring deliberately because SQL text, identities, dependencies,
and sizes can reveal operational information. An application does not need
monitoring privileges to execute a cacheable query.

\subsection{A controlled evaluation and cleanup}

\begin{enumerate}
  \item Check the uncached result and capture its scan or join plan with
    caching disabled.
  \item Use exact mode for the current BioSeqDB adapter. Run cold and warm
    executions of the same request, and read their actual cache actions.
  \item Compare production, hit, wait, reuse, and fallback counter changes.
    Check the entries' state, completion time, retention duration, rows,
    and payload sizes.
  \item Make a representative write and compare the next cached result
    with an uncached execution on the changed data.
  \item With a sufficiently long query, cancel a waiting test session and
    check that the producer can continue. Inspect failure and cleanup state.
  \item Keep manual invalidation restricted to a superuser. Leave caching
    disabled for one-off results unlikely to be consumed before expiry or
    invalidation.
\end{enumerate}

The table remains available throughout this tutorial, including the concurrent
session and write examples. Remove its temporary setup when those checks end.

\par\noindent\begin{minipage}{\linewidth}
\begin{lstlisting}[style=bioseq-appendix-sql,
  caption={Remove the example objects and reset session settings.}]
DROP TABLE cache_stats_before;
DROP TABLE cache_reads;
RESET bioseq_cache.query_cache.mode;
RESET bioseq_cache.query_cache.completed_result_ttl_ms;
RESET bioseq.index.insert_maintenance.enabled;
\end{lstlisting}
\end{minipage}
